\documentclass[aspectratio=169,11pt]{beamer}

% ============================================================
% R3. Reunião de 23/10/2026: ciclo v8 do Qwen3-4B (19 a 23/10/2026).
% Fontes dos números (todas em /home/nees/Documents/VSCodigo/TreinamentoModeloQuestoes):
%   ciclo_v8/SPEC.md                                        (decisões, auditoria, resultados)
%   outputs/melhoria_v8/F/apuracao_parcial.json             (professor, 222 itens do 4B v7b)
%   outputs/melhoria_v8/validacao_jev_fora_da_amostra.json  (calibração do Jev)
%   outputs/melhoria_v8/calibracao_maritaca.json            (calibração da Maritaca)
%   outputs/melhoria_v8/avaliar_jev_*.jsonl + cache Maritaca (juízes pareados, 800/400 pares)
%   outputs/relatorios/VEREDITO_v8rftb_vs_v7b_multiseed.json (gate, 90 pares)
%   outputs/relatorios/VEREDITO_v8rftb_vs_v7b_multiseed_G1_todos_reprovado.json
%   ciclo_v7/SPEC.md, outputs/gemma3/amostra_humana_s44/apuracao.json (modelos antigos)
%   data/aval_{v7b,v8rft,v8rftb}.jsonl                      (letras e "Nenhuma", 800 amostras)
% ============================================================

\usetheme{Madrid}
\usecolortheme{default}
\setbeamertemplate{navigation symbols}{}
\setbeamertemplate{footline}[frame number]

\usepackage[utf8]{inputenc}
\usepackage[T1]{fontenc}
\usepackage[brazil]{babel}
\usepackage{booktabs}
\usepackage{tabularx}
\usepackage{tikz}
\usetikzlibrary{arrows.meta, positioning, shapes.geometric, fit, backgrounds, calc}
\usepackage{pgfplots}
\pgfplotsset{compat=1.17}
\pgfkeys{/pgf/number format/.cd, use comma, 1000 sep={}}
\usepackage{fontawesome5}
\usepackage{graphicx}
\usepackage{xcolor}
\usepackage{array}
\usepackage{colortbl}
\usepackage{tcolorbox}
\tcbuselibrary{skins}

% ---------- paleta (mesma do R1 e do R2) ----------
\definecolor{okgreen}{RGB}{34,139,34}
\definecolor{warnorange}{RGB}{204,102,0}
\definecolor{badred}{RGB}{178,34,34}
\definecolor{basegray}{RGB}{90,90,90}
\definecolor{accentblue}{RGB}{31,78,121}
\definecolor{lightblue}{RGB}{221,235,247}
\definecolor{lightgreen}{RGB}{226,240,217}
\definecolor{lightred}{RGB}{248,225,225}
\definecolor{lightorange}{RGB}{252,235,214}

\setbeamercolor{structure}{fg=accentblue}

\newcommand{\okmark}{{\color{okgreen}\faCheckCircle}}
\newcommand{\warnmark}{{\color{warnorange}\faExclamationTriangle}}
\newcommand{\badmark}{{\color{badred}\faTimesCircle}}
\newcommand{\pendmark}{{\color{basegray}\faHourglassHalf}}

% Ilustrações externas: se o arquivo ainda não existir em img/, mostra um
% quadro tracejado com o nome do arquivo a baixar (Freepik/Pixabay/Pexels).
\newcommand{\figext}[3][0.3\textwidth]{%
  \IfFileExists{img/#2}{\includegraphics[width=#1]{img/#2}}{%
    \begin{tikzpicture}
      \node[draw=basegray, dashed, rounded corners, minimum width=#1, minimum height=0.55*#1,
            align=center, text=basegray, font=\tiny, inner sep=3pt]
        {\faImage\\[2pt]\texttt{\detokenize{img/#2}}\\[1pt]#3};
    \end{tikzpicture}}}

\newtcolorbox{caixa}[2][]{enhanced, colback=#2!40!white, colframe=#2!70!black, boxrule=0.5pt,
  arc=2pt, left=4pt, right=4pt, top=3pt, bottom=3pt, fontupper=\small, #1}

\tikzset{
  etapa/.style={draw=accentblue, fill=lightblue, rounded corners, align=center,
                font=\scriptsize, minimum height=0.9cm, text width=2.2cm},
  seta/.style={-{Stealth[length=2mm]}, thick, accentblue},
}

\title[R3, 23/10/2026]{R3. Reunião de acompanhamento}
\subtitle{Ciclo v8 do Qwen3-4B: base de dados, juízes, gate e expectativa para o aplicativo}
\author{Relatório de execução de 19 a 23 de outubro de 2026}
\institute{Qwen3-4B + QLoRA + GGUF, geração offline de questões SAEB de Matemática}
\date{23 de outubro de 2026}

\begin{document}

% ------------------------------------------------------------
\begin{frame}
  \titlepage
\end{frame}

% ------------------------------------------------------------
\begin{frame}{Roteiro}
  \begin{columns}[T]
    \begin{column}{0.55\textwidth}
      \begin{enumerate}
        \item Resumo executivo
        \item Ponto de partida: o que o professor mediu
        \item Como decidimos o caminho
        \item Melhorias na base de dados
        \item Juízes automáticos calibrados
        \item Resultados: juízes pareados e gate
        \item Comparação com os modelos anteriores
        \item Expectativa para o aplicativo
        \item Riscos e próximos passos
      \end{enumerate}
    \end{column}
    \begin{column}{0.4\textwidth}
      \centering
      \figext[0.95\textwidth]{reuniao_equipe_relatorio.png}{reunião de equipe com gráficos}
    \end{column}
  \end{columns}
\end{frame}

% ------------------------------------------------------------
\section{Resumo}
\begin{frame}{Resumo executivo}
  \begin{columns}[T]
    \begin{column}{0.32\textwidth}
      \begin{caixa}{lightblue}
        \textbf{\faDatabase\ Base}\\[2pt]
        2.918 questões reordenadas para \emph{resolver antes de escrever as alternativas}, mais
        875 questões novas aprovadas por juiz. Total de 3.793, com gabarito em 20\% por letra.
      \end{caixa}
    \end{column}
    \begin{column}{0.32\textwidth}
      \begin{caixa}{lightgreen}
        \textbf{\faBalanceScale\ Juízes}\\[2pt]
        Jev: 64,6\% $\rightarrow$ \textbf{73,4\%} de aceitação (800 pares, $p<0{,}001$).
        Maritaca: 60,2\% $\rightarrow$ 66,5\% (400 pares, $p=0{,}065$). Os dois apontam a mesma direção.
      \end{caixa}
    \end{column}
    \begin{column}{0.32\textwidth}
      \begin{caixa}{lightgreen}
        \textbf{\faCheckCircle\ Gate}\\[2pt]
        v8rftb contra v7b: \textbf{PROMOVIDO}. Consistência sobe de 93,6\% para 100\% e as
        falhas caem de 3 para 1 em 90.
      \end{caixa}
    \end{column}
  \end{columns}
  \vspace{4pt}
  \begin{caixa}{lightred}
    \textbf{\faExclamationTriangle\ O professor não confirmou o ganho:} núcleo de \textbf{65,4\%} (51 de 78,
    IC95 [54,3; 75,0]), contra 68,9\% do v7b. A diferença está dentro do ruído e a meta de 80\% não foi atingida.
    \textbf{Recomendação: não publicar o v8rftb.} O 9º ano e a Geometria seguem como pontos fracos.
  \end{caixa}
\end{frame}

% ------------------------------------------------------------
\section{Ponto de partida}
\begin{frame}{Ponto de partida: 222 questões do 4B v7b avaliadas pelo professor}
  \begin{columns}[T]
    \begin{column}{0.48\textwidth}
      \small
      O professor avaliou às cegas 222 questões do melhor modelo da época (Qwen3-4B v7b).
      A questão é válida no critério \textbf{núcleo} quando a matemática está correta, só uma
      alternativa responde e os dados bastam.
      \vspace{4pt}
      \begin{itemize}
        \item Núcleo: \textbf{68,9\%}, IC95 [62,5; 74,6]
        \item A base de treino tem 97,9\% de núcleo. O gargalo é o modelo, não os dados.
        \item Hiperparâmetros não mostraram efeito acima do ruído entre seeds.
      \end{itemize}
    \end{column}
    \begin{column}{0.5\textwidth}
      \centering
      \begin{tikzpicture}
        \begin{axis}[
          xbar, width=0.98\textwidth, height=4.6cm,
          xmin=0, xmax=100, xlabel={\scriptsize validade núcleo (\%)},
          symbolic y coords={9º ano, 5º ano, 2º ano, Geometria, Álgebra, Números},
          ytick=data, yticklabel style={font=\scriptsize}, xticklabel style={font=\scriptsize},
          nodes near coords, nodes near coords style={font=\tiny},
          bar width=7pt, enlarge y limits=0.12,
        ]
          \addplot[fill=accentblue!70, draw=accentblue] coordinates {
            (54,9º ano) (70,5º ano) (82,2º ano) (42,Geometria) (51,Álgebra) (58,Números)};
        \end{axis}
      \end{tikzpicture}

      {\scriptsize 9º ano e Geometria concentram as falhas.}
    \end{column}
  \end{columns}
\end{frame}

% ------------------------------------------------------------
\begin{frame}{Por que as 69 questões falharam}
  \begin{columns}[c]
    \begin{column}{0.55\textwidth}
      \centering
      \begin{tikzpicture}
        \begin{axis}[
          xbar, width=0.62\textwidth, height=4.4cm,
          xmin=0, xmax=30, xlabel={\scriptsize questões},
          symbolic y coords={Valor fora das alternativas, Mais de uma correta, Conta ou conceito errado, Enunciado mal posto},
          ytick=data, yticklabel style={font=\scriptsize, align=right, text width=3.2cm},
          xticklabel style={font=\scriptsize},
          nodes near coords, nodes near coords style={font=\scriptsize},
          bar width=9pt, enlarge y limits=0.18,
        ]
          \addplot[fill=badred!60, draw=badred] coordinates {
            (12,Valor fora das alternativas) (14,Mais de uma correta)
            (19,Conta ou conceito errado) (24,Enunciado mal posto)};
        \end{axis}
      \end{tikzpicture}
    \end{column}
    \begin{column}{0.42\textwidth}
      \small
      A causa mais comum é o enunciado mal posto, com dados contraditórios ou inventados.
      Depois vêm a conta errada e a resposta duplicada.
      \vspace{4pt}

      Reordenar o JSON resolve só a última barra, que vale cerca de \textbf{5 pontos}.
      Para as outras três, era preciso filtrar o que o modelo aprende. Daí veio o RFT (\emph{Rejection Fine-Tuning}, ou ajuste fino por rejeição): o modelo gera várias questões, um juiz descarta as ruins e o modelo é treinado de novo só com as aprovadas.
    \end{column}
  \end{columns}
\end{frame}

% ------------------------------------------------------------
\section{Decisão}
\begin{frame}{Como decidimos o caminho}
  \centering
  \begin{tikzpicture}[node distance=0.55cm and 0.5cm]
    \node[etapa] (f) {\textbf{F. Medir}\\professor, 222 itens};
    \node[etapa, right=of f] (j) {\textbf{Jev}\\consultado 2 vezes};
    \node[etapa, right=of j, fill=lightgreen, draw=okgreen] (b) {\textbf{B. RFT}\\filtrar o que o modelo aprende};
    \node[etapa, below=of b, fill=lightgreen, draw=okgreen] (a) {\textbf{A. Ordem v8}\\resolver antes das alternativas};
    \node[etapa, right=of b] (t) {\textbf{v8rftb}\\um único treino};
    \draw[seta] (f) -- (j);
    \draw[seta] (j) -- node[above, font=\tiny] {1º} (b);
    \draw[seta] (j) |- node[pos=0.75, below, font=\tiny] {junto} (a);
    \draw[seta] (b) -- (t);
    \draw[seta] (a) -| (t);
  \end{tikzpicture}
  \vspace{6pt}

  \small
  \begin{tabularx}{0.92\textwidth}{@{}lX@{}}
    \toprule
    Consulta & Resultado \\
    \midrule
    1ª, antes da medição & A primeiro (probabilidade média 0,76). A medição F era pré-requisito (0,94). \\
    2ª, com a medição    & B primeiro nas 4 ordens testadas (B 0,67 contra A 0,30). A e B se complementam. \\
    Hiperparâmetros      & Sem evidência de ganho acima do ruído entre seeds (0,05). Ficaram de fora. \\
    \bottomrule
  \end{tabularx}
\end{frame}

% ------------------------------------------------------------
\section{Base de dados}
\begin{frame}{Melhorias na base de dados}
  \centering
  \begin{tikzpicture}[node distance=0.6cm]
    \tikzset{etapa/.append style={minimum height=0.75cm, font=\tiny}}
    \node[etapa, text width=2.3cm] (v7) {\textbf{v7b\_fix1}\\2.948 questões, núcleo 97,9\%};
    \node[etapa, right=of v7, text width=2.5cm] (v8) {\textbf{v8}\\2.918 questões, ordem nova};
    \node[etapa, right=of v8, text width=2.5cm, fill=lightorange, draw=warnorange] (rft) {\textbf{RFT}\\+875 aprovadas pelo Jev};
    \node[etapa, right=of rft, text width=2.5cm, fill=lightgreen, draw=okgreen] (bal) {\textbf{v8\_rft\_bal}\\3.793 questões, 20\% por letra};
    \draw[seta] (v7) -- (v8);
    \draw[seta] (v8) -- (rft);
    \draw[seta] (rft) -- (bal);
  \end{tikzpicture}
  \vspace{2pt}

  \scriptsize
  \begin{tabularx}{\textwidth}{@{}lX@{}}
    \toprule
    Mudança & Por que importa \\
    \midrule
    Resolução antes das alternativas & O valor calculado já existe quando o modelo escreve as opções. \\
    Uma questão por exemplo (47 divididos) & É o formato real do aplicativo, que pede uma questão por chamada. \\
    Resolução sem citar ``alternativa B'' & As alternativas ainda não existem nessa ordem. Isso também permite embaralhar as opções no app. \\
    Trava do verificador & Se a limpeza mudasse o veredito, a questão saía. Saíram 30 injetadas e nenhum item real do banco. \\
    Extras com letras permutadas & Os extras chegaram com 34,5\% de gabarito A. A permutação corrige para 20\% cada. \\
    \bottomrule
  \end{tabularx}
\end{frame}

% ------------------------------------------------------------
\begin{frame}{RFT (ajuste fino por rejeição): o modelo gera, o juiz filtra, o modelo reaprende}
  \centering
  \begin{tikzpicture}
    \pgfmathsetmacro{\escala}{7.5/2400}
    \foreach \y/\n/\rot/\cor in {
        0/2400/geradas pelo 4B v7b (400 prompts $\times$ 6)/accentblue,
        -0.75/2067/passaram no filtro determinístico/accentblue,
        -1.5/1461/julgadas pelo Jev/warnorange,
        -2.25/1109/aceitas pelo Jev/okgreen,
        -3.0/875/entraram no treino (teto de 30\% da base)/okgreen} {
      \fill[\cor!55, draw=\cor] (0,\y) rectangle ({\n*\escala},{\y+0.55});
      \node[anchor=west, font=\scriptsize, text=white] at (0.1,{\y+0.275}) {\textbf{\n}};
      \node[anchor=west, font=\scriptsize] at ({\n*\escala+0.15},{\y+0.275}) {\rot};
    }
  \end{tikzpicture}
  \vspace{6pt}

  \small
  \begin{itemize}
    \item 606 amostras ficaram sem julgamento quando o crédito do OpenRouter acabou. Mesmo assim, as 79 habilidades receberam extras.
    \item Também saíram 282 pares bom/ruim, guardados para um treino ORPO futuro.
  \end{itemize}
\end{frame}

% ------------------------------------------------------------
\section{Juízes}
\begin{frame}{Dois juízes calibrados contra o professor}
  \small
  Antes de usar um juiz para decidir, conferimos se ele concorda com o professor nas mesmas 222 questões.
  \vspace{6pt}

  \centering
  \begin{tabular}{@{}lcc@{}}
    \toprule
    & \textbf{Jev} (TypeSafe) & \textbf{Maritaca} (sabiá-4) \\
    \midrule
    Concordância com o professor ($\kappa$) & \textbf{0,795} & 0,725 \\
    Inválidos que o juiz pega (recall)      & 81,2\% & \textbf{90,9\%} \\
    Válidos que o juiz perde                & \textbf{3,9\%} & 13,7\% \\
    Válidos entre os aceitos                & 91,9\% & \textbf{95,7\%} \\
    Custo da calibração                     & US\$ 0,01 & R\$ 5,36 \\
    Papel neste ciclo                       & filtro do RFT e 1º juiz & 2º juiz, independente do filtro \\
    \bottomrule
  \end{tabular}
  \vspace{6pt}

  \raggedright
  \small A meta era $\kappa \geq 0{,}6$, e os dois passaram. A Maritaca é mais rígida e não participou do filtro
  de treino, então ela confirma o ganho sem viés de construção. Os dois rodam por API, por isso servem para medir e
  filtrar dados, mas não vão para o celular.
\end{frame}

% ------------------------------------------------------------
\section{Resultados}
\begin{frame}{Resultados: os mesmos prompts e seeds nos dois modelos}
  \begin{columns}[c]
    \begin{column}{0.55\textwidth}
      \centering
      \begin{tikzpicture}
        \begin{axis}[
          ybar, width=\textwidth, height=5.2cm,
          ymin=0, ymax=85, ylabel={\scriptsize aceitação (\%)},
          symbolic x coords={Jev, Maritaca},
          xtick=data, xticklabel style={font=\scriptsize}, yticklabel style={font=\scriptsize},
          nodes near coords, nodes near coords style={font=\tiny},
          bar width=16pt, enlarge x limits=0.45,
          legend style={font=\scriptsize, at={(0.5,1.02)}, anchor=south, legend columns=3},
        ]
          \addplot[fill=basegray!50, draw=basegray] coordinates {(Jev,64.6) (Maritaca,60.2)};
          \addplot[fill=warnorange!60, draw=warnorange] coordinates {(Jev,73.6)};
          \addplot[fill=okgreen!60, draw=okgreen] coordinates {(Jev,73.4) (Maritaca,66.5)};
          \legend{v7b, v8rft, v8rftb}
        \end{axis}
      \end{tikzpicture}
    \end{column}
    \begin{column}{0.43\textwidth}
      \footnotesize
      \begin{tabular}{@{}lcc@{}}
        \toprule
        Pareado & Jev & Maritaca \\
        \midrule
        Pares & 800 & 400 \\
        Só o v7b aceita & 110 & 72 \\
        Só o v8rftb aceita & \textbf{180} & \textbf{97} \\
        McNemar $p$ & $<0{,}001$ & 0,065 \\
        \bottomrule
      \end{tabular}
      \vspace{6pt}

      \small
      O balanceamento não custou qualidade, porque o v8rft e o v8rftb ficam empatados no Jev.
      No 9º ano, a aceitação do Jev sobe de 55\% para 68\%.
    \end{column}
  \end{columns}
\end{frame}

% ------------------------------------------------------------
\begin{frame}{Gate mecânico: v8rftb contra v7b (3 rodadas, 90 pares)}
  \centering
  \footnotesize
  \begin{tabularx}{\textwidth}{@{}l>{\raggedright\arraybackslash}Xccc@{}}
    \toprule
    Gate & O que verifica & v7b & v8rftb & \\
    \midrule
    G1* & Estrutura do JSON nas questões entregues & 100\% & 100\% (n=89) & \okmark \\
    G2* & Falhas depois das tentativas (teto 5\%) & 3/90 & \textbf{1/90} & \okmark \\
    G3* & Gabarito consistente com a conta & 93,6\% & \textbf{100\%} & \okmark \\
    G4  & Dependência de figura & 0 & 0 & \okmark \\
    G6* & Letra mais frequente (limite: base + 5 pontos) & 40,0\% & 41,1\% & \okmark \\
    G7* & Sem esquecimento no 5º e 9º anos & 100\% & 100\% & \okmark \\
    \midrule
    G5  & Dificuldade (informativo) & 100\% & 100\% & \okmark \\
    G8  & Velocidade (informativo) & 13,8 tok/s & 15,0 tok/s & \okmark \\
    \bottomrule
  \end{tabularx}
  \vspace{6pt}

  \begin{caixa}{lightgreen}
    \textbf{Veredito: PROMOVIDO.} O gate garante que nada piorou, mas não mede se a matemática está certa.
    Quem mede isso é o professor.
  \end{caixa}
\end{frame}

% ------------------------------------------------------------
\begin{frame}{Mudanças nas regras do gate, com transparência}
  \small
  \begin{tabularx}{\textwidth}{@{}lX>{\raggedright\arraybackslash}p{3.3cm}@{}}
    \toprule
    Data & Mudança e motivo & Tipo \\
    \midrule
    22/10 & G5 e G8 viram informativos e G9 sai. G5 só copiava o rótulo pedido (100\% nos dois lados).
            G8 reprovou 6 de 18 vereditos por tamanho do modelo ou aquecimento da GPU. & Decisão do responsável, antes de medir o candidato \\
    23/10 & O G1* passa a contar só as questões entregues. O único item reprovado tinha B igual a D, mas já tinha sido
            descartado como falha e contado no G2*. Com 90 pares, o limite de 99\% virava tolerância zero. & Revisão depois de uma reprovação, documentada no código \\
    \bottomrule
  \end{tabularx}
  \vspace{6pt}

  \begin{itemize}\small
    \item Uma questão entregue com estrutura ruim continua reprovando. Três testes novos garantem isso.
    \item Com \texttt{--gates-completos}, a regra antiga volta e o veredito antigo (NÃO PROMOVIDO) se reproduz.
    \item O veredito reprovado ficou guardado como evidência.
  \end{itemize}
\end{frame}

% ------------------------------------------------------------
\section{Comparação}
\begin{frame}{Comparação com os modelos anteriores}
  \centering
  \scriptsize
  \begin{tabularx}{\textwidth}{@{}l c c c >{\raggedright\arraybackslash}X@{}}
    \toprule
    Modelo & Tamanho & Núcleo (professor) & Gate & Melhor contribuição \\
    \midrule
    Gemma 3 1B zero-shot & 0,8 GB & 15\% (n=20) & base & Pipeline Gemma pronto e trava que impede comparar famílias diferentes \\
    Gemma 3 1B fine-tune & 0,8 GB & 0\% (n=20) & \badmark\ 3 seeds & Mostrou que 1B não comporta a tarefa \\
    Qwen3-1.7B v3 & 1,1 GB & 50\% (n=20)\textsuperscript{a} & \okmark & Primeiro modelo publicado, gramática GBNF e gate multiseed \\
    Qwen3-1.7B v6 & 1,1 GB & 35\% (n=20) & \okmark & Mostrou que passar no gate não garante qualidade, e a amostra humana passou a decidir \\
    Qwen3-1.7B v7b & 1,1 GB & 35 a 50\% (n=20) & \okmark & Base com 97,9\% de núcleo, o que mostrou que o gargalo é o modelo \\
    Qwen3-4B zero-shot & 2,4 GB & 65 a 70\% (n=20) & base & Teste de teto: o tamanho do modelo pesa mais que a curadoria \\
    Qwen3-4B v7b & 2,5 GB & \textbf{68,9\%} (n=222) & \okmark & Primeira medida confiável, taxonomia de falhas e calibração dos juízes \\
    \rowcolor{lightgreen}
    Qwen3-4B v8rftb & 2,5 GB & 65,4\% (n=78) & \okmark & Gate e juízes melhoraram, mas o professor não viu ganho. Mostrou o limite do filtro por juiz de API \\
    \bottomrule
  \end{tabularx}
  \vspace{4pt}

  \raggedright
  \tiny \textsuperscript{a} Valor do baseline\_v1 do 1.7B na amostra contra o v6. Com n=20, o intervalo de confiança tem cerca de $\pm$20 pontos,
  por isso as linhas do 1.7B se sobrepõem. Tamanho do 4B fine-tunado: GGUF Q4\_K\_M de 2,4 GB mais adaptador LoRA de 66 MB.
\end{frame}

% ------------------------------------------------------------
\begin{frame}{A evolução em um gráfico}
  \centering
  \begin{tikzpicture}
    \begin{axis}[
      ybar, width=0.92\textwidth, height=5.6cm,
      ymin=0, ymax=100, ylabel={\scriptsize validade núcleo (\%)},
      symbolic x coords={Gemma 1B FT, 1.7B v6, 1.7B v7b, 4B ZS, 4B v7b, 4B v8rftb}, xticklabel style={font=\scriptsize}, yticklabel style={font=\scriptsize},
      nodes near coords, nodes near coords style={font=\tiny},
      bar width=20pt, enlarge x limits=0.1, bar shift=0pt,
      xtick={Gemma 1B FT, 1.7B v6, 1.7B v7b, 4B ZS, 4B v7b, 4B v8rftb},
    ]
      \addplot[fill=accentblue!60, draw=accentblue] coordinates {
        (Gemma 1B FT,0) (1.7B v6,35) (1.7B v7b,35) (4B ZS,65) (4B v7b,68.9)};
      \addplot[fill=warnorange!50, draw=warnorange] coordinates {(4B v8rftb,65.4)};
      \draw[badred, dashed, thick] (axis cs:Gemma 1B FT,80) -- (axis cs:4B v8rftb,80)
        node[pos=0.03, above, font=\tiny, text=badred] {meta 80\%};
    \end{axis}
  \end{tikzpicture}

  \footnotesize Todas as barras são medidas do professor (n=20 nos modelos pequenos, n=222 no 4B v7b e n=78 no v8rftb).
  O v8rftb (65,4\%) e o v7b (68,9\%) têm intervalos que se sobrepõem, e nenhum chega à meta de 80\%.
\end{frame}

% ------------------------------------------------------------
\section{Aplicativo}
\begin{frame}{Expectativa para o aplicativo mobile}
  \begin{columns}[T]
    \begin{column}{0.64\textwidth}
      \scriptsize
      \begin{tabular}{@{}lll@{}}
        \toprule
        Item & Qwen3-1.7B (hoje) & Qwen3-4B v8rftb \\
        \midrule
        Arquivo do modelo & 1,1 GB & 2,4 GB + 66 MB \\
        Memória estimada\textsuperscript{b} & cerca de 1,4 GB & cerca de 3 GB \\
        Velocidade (CPU do PC) & 27 tok/s & 15 tok/s \\
        Tempo por questão (PC) & cerca de 8 s & cerca de 21 s \\
        Núcleo (professor) & 35 a 50\% & 65\% (n=78) \\
        Gramática enviada & \texttt{questao.gbnf} & \texttt{questao\_v8.gbnf} \\
        \bottomrule
      \end{tabular}
      \vspace{2pt}

      \tiny \textsuperscript{b} Pesos mais o cache de 2.048 tokens (cerca de 0,3 GB no 4B). Estimativa de cálculo, ainda sem medida no aparelho.
    \end{column}
    \begin{column}{0.33\textwidth}
      \centering
      \figext[0.9\textwidth]{celular_offline_questoes.png}{celular com questão de matemática, modo avião}
    \end{column}
  \end{columns}
  \vspace{4pt}

  \begin{caixa}{lightblue}
    \footnotesize
    \textbf{O que esperar:} correção parecida com a do v7b (65 a 69\%), bem acima do 1.7B, ao custo de cerca de 2,5 vezes mais tempo por questão e 2 vezes mais memória.
    Pela estimativa, o 4B pede aparelhos com 6 GB de RAM ou mais. Em celulares, a velocidade deve ficar abaixo da do PC.
    Por isso é preciso medir em um aparelho real antes de publicar.
  \end{caixa}
\end{frame}

% ------------------------------------------------------------
\begin{frame}{O que muda na integração com o app}
  \small
  \begin{enumerate}
    \item \textbf{Trocar a gramática junto com o modelo.} O v8 só funciona com \texttt{questao\_v8.gbnf}, que o app envia inline.
          As chaves do JSON continuam as mesmas, então a leitura por nome de chave não muda.
    \item \textbf{Embaralhar as alternativas no app.} A resolução do v8 não cita letras, então o app pode sortear a
          ordem das opções e remapear o gabarito. Assim o viés de letra deixa de importar para o aluno.
    \item \textbf{Descartar o que vier como falha.} O pipeline já marca e descarta. O aluno nunca vê essa questão.
    \item \textbf{Escolher o empacotamento.} Uma opção é usar o modelo base e o adaptador separados, como foi medido. A outra é um GGUF único
          com o adaptador mesclado, que exige passar pelo gate de novo.
    \item \textbf{Manter o 1.7B como alternativa} para aparelhos com pouca memória.
  \end{enumerate}
\end{frame}

% ------------------------------------------------------------
\section{Riscos e próximos passos}
\begin{frame}{Pontos de atenção}
  \begin{columns}[T]
    \begin{column}{0.5\textwidth}
      \centering
      \begin{tikzpicture}
        \begin{axis}[
          ybar, width=\textwidth, height=4.6cm,
          ymin=0, ymax=55, ylabel={\scriptsize \% das questões},
          symbolic x coords={Gabarito A, Nenhuma das anteriores},
          xtick=data, xticklabel style={font=\scriptsize}, yticklabel style={font=\scriptsize},
          nodes near coords, nodes near coords style={font=\tiny},
          bar width=11pt, enlarge x limits=0.45,
          legend style={font=\tiny, at={(0.5,1.02)}, anchor=south, legend columns=3},
        ]
          \addplot[fill=basegray!50, draw=basegray] coordinates {(Gabarito A,33.7) (Nenhuma das anteriores,14.2)};
          \addplot[fill=warnorange!60, draw=warnorange] coordinates {(Gabarito A,44.0) (Nenhuma das anteriores,26.9)};
          \addplot[fill=okgreen!60, draw=okgreen] coordinates {(Gabarito A,41.8) (Nenhuma das anteriores,28.3)};
          \legend{v7b, v8rft, v8rftb}
        \end{axis}
      \end{tikzpicture}

      {\tiny 800 amostras pareadas por modelo. O equilíbrio seria 20\% por letra.}
    \end{column}
    \begin{column}{0.48\textwidth}
      \small
      \begin{itemize}
        \item \warnmark\ \textbf{O viés de letra persiste.} Balancear o treino levou o gabarito A de 44,0\% para só 41,8\%.
              O gate passou porque o v7b também deu 40\% nas rodadas dele. Embaralhar as opções no app resolve isso.
        \item \warnmark\ \textbf{Nenhuma das anteriores} dobrou. O corpus tem cerca de 21\%, então o modelo está copiando o treino.
              O professor deve dizer se isso prejudica.
        \item \warnmark\ \textbf{9º ano e Geometria} continuam como os pontos fracos.
      \end{itemize}
    \end{column}
  \end{columns}
\end{frame}

% ------------------------------------------------------------
\begin{frame}{Resultado do professor e diagnóstico das 27 reprovadas}
  \begin{columns}[T]
    \begin{column}{0.5\textwidth}
      \footnotesize
      \begin{tabular}{@{}lcc@{}}
        \toprule
        Recorte & Núcleo & n \\
        \midrule
        Geral & \textbf{65,4\%} & 78 \\
        2º ano & 83\% & 23 \\
        5º ano & 59\% & 22 \\
        9º ano & 50\% & 26 \\
        Geometria & 46\% & 13 \\
        Grandezas e medidas & 46\% & 13 \\
        Números & 58\% & 33 \\
        Verificador \texttt{ok} & 79\% & 42 \\
        Verificador não verificável & 50\% & 36 \\
        \bottomrule
      \end{tabular}
    \end{column}
    \begin{column}{0.48\textwidth}
      \small
      Causas das 27 reprovadas, no texto do professor:
      \begin{itemize}\footnotesize
        \item cerca de 11 com enunciado contraditório (passageiros, troco, totais),
        \item cerca de 6 com mais de uma resposta correta,
        \item cerca de 4 com conta ou resolução errada e
        \item o restante com situação mal definida.
      \end{itemize}
      \footnotesize
      \textbf{Nenhuma das anteriores} não pesou (59\% contra 68\%, dentro do ruído).
      O filtro determinístico atual só marca o que não consegue verificar, então não separa bem.
    \end{column}
  \end{columns}
\end{frame}

% ------------------------------------------------------------
\begin{frame}{Próximos passos}
  \centering
  \begin{tikzpicture}[node distance=0.5cm and 0.45cm]
    \node[etapa, fill=lightred, draw=badred] (r) {\textbf{v8rftb}\\65,4\%\\abaixo de 80\%};
    \node[etapa, right=of r, fill=lightgreen, draw=okgreen, text width=2.6cm] (c) {\textbf{Auditoria dos extras}\\feita: 88,3\% de núcleo};
    \node[etapa, right=of c, fill=lightorange, draw=warnorange, text width=2.6cm] (j) {\textbf{Gerar e conferir}\\resolver às cegas, no app};
    \node[etapa, right=of j, text width=2.6cm] (o) {\textbf{ORPO}\\descartado por ora};
    \draw[seta] (r) -- (c);
    \draw[seta] (c) -- (j);
    \draw[seta] (j) -- (o);
  \end{tikzpicture}
  \vspace{4pt}

  \footnotesize
  \begin{tabularx}{\textwidth}{@{}>{\raggedright\arraybackslash}p{5.6cm}X@{}}
    \toprule
    Ação & Motivo \\
    \midrule
    Manter o v7b como referência & Mesmo núcleo do v8rftb, sem adotar um modelo não confirmado \\
    Auditoria dos extras do RFT, já feita & Núcleo de 88,3\% (n=60) contra 65,4\% do modelo: os dados são bons e o gargalo é a geração \\
    Validar a resposta na hora de gerar, dentro do app & Gerar, resolver de novo às cegas e aceitar só se concordar, sem API \\
    Medir em um celular real e embaralhar as alternativas & Pendências do app, independentes do modelo \\
    Descartar o ORPO por ora & Os dados já são bons e os pares vêm do mesmo filtro, então não atacam a causa \\
    \bottomrule
  \end{tabularx}
\end{frame}

% ------------------------------------------------------------
\appendix
\begin{frame}{Apêndice: arquivos de evidência}
  \scriptsize
  \begin{tabularx}{\textwidth}{@{}>{\ttfamily}lX@{}}
    \toprule
    Arquivo & Conteúdo \\
    \midrule
    ciclo\_v8/SPEC.md & Hipótese, decisões, auditoria e resultados do ciclo \\
    outputs/melhoria\_v8/F/ & Folha do professor (222 itens) e apuração \\
    outputs/melhoria\_v8/validacao\_jev\_fora\_da\_amostra.json & Calibração do Jev \\
    outputs/melhoria\_v8/calibracao\_maritaca.json & Calibração da Maritaca \\
    data/train\_curado\_v8\_rft\_bal.jsonl & Corpus treinado (3.793 questões) \\
    data/aval\_\{v7b,v8rftb\}.jsonl & 800 amostras pareadas por modelo \\
    outputs/relatorios/VEREDITO\_v8rftb\_vs\_v7b\_multiseed.json & Gate PROMOVIDO \\
    outputs/relatorios/VEREDITO\_v8rftb\_\dots\_G1\_todos\_reprovado.json & Veredito anterior à revisão do G1* \\
    outputs/qwen3\_4b/lora\_4b\_v8rftb\_s44.f16.gguf & Adaptador do candidato (66 MB) \\
    \bottomrule
  \end{tabularx}
\end{frame}

\end{document}
